\clearpage\begingroup\SourceLeaf{450b}\ArithmeticHead{119}{R}{DIE ALLGEMEINE LEHRE VON DEN KRUMMEN FLÄCHEN.}
\begin{center}\textbf{34.}\par\smallskip
\textls[100]{Bibliographischer Anhang.}\end{center}
\begingroup\small
\noindent\hspace*{\parindent}Die von \textsc{Gauss} Ende 1822 an die Kopenhagener Societät der Wissenschaften eingesandte Abhandlung
über die Abbildung krummer Flächen hatte zwar den Preis erhalten, allein die Gesellschaft überliess es
den Preisträgern, für die Veröffentlichung zu sorgen, und so ist die Preisschrift erst 1825 im dritten und
letzten Heft der von \textsc{Schumacher} als Ergänzung der Astronomischen Nachrichten herausgegebenen Astronomischen
Abhandlungen erschienen. Sie ist abgedruckt in den Werken, Bd. IV, 1873, 2. Abdruck 1880,
S. 189—216. Eine Übersetzung ins Englische, wahrscheinlich von \textsc{Francis Baily} (1764—1844), ist 1828
erschienen:
\par
General solution of the problem: to represent the parts of a given surface on another given surface,
so that the smallest parts of the representation shall be similar to the corresponding parts of the surface
represented. By \textsc{C. F. Gauss}. Answer to the Prize Question proposed by the Royal Society of sciences
at Copenhagen. The philosophical magazine, new series, vol. 4, London 1828, S. 104—113, 206—215.
\par
Im Jahre 1894 ist die Abhandlung von \textsc{A. Wangerin} neu herausgegeben worden; sie findet sich im
Hefte 55 von \textsc{Ostwalds} Klassikern der exakten Wissenschaften: Über Kartenprojection, Abhandlungen von
\textsc{Lagrange} (1779) und \textsc{Gauss} (1822), Leipzig 1894, S. 57—81.
\par
Die Abhandlung über die allgemeine Lehre von den krummen Flächen hat \textsc{Gauss} am 8. Oktober
1827 der Göttinger Societät vorgelegt. Eine von \textsc{Gauss} selbst verfasste Anzeige erschien am 5. November
1827 in den Göttingischen Gelehrten Anzeigen, Stück 177, S. 1761—1768; sie ist abgedruckt in den Werken,
Bd. IV, S. 341—347. Eine Übersetzung der Selbstanzeige ist schon 1829 von \textsc{Francis Baily} herausgegeben
worden:
\par
Account of a paper by Prof. \textsc{Gauss}, intitled: Disquisitiones generales circa superficies curvas, communicated
to the Royal Society of Göttingen on the 8th of october 1827, The philosophical magazine, new
series, vol. 3, London 1828, S. 331—336.
\par
Man vgl. hierzu den Brief von \textsc{Olbers} an \textsc{Gauss} vom 2. Juli 1828 und dessen Antwort Ende Juli
1828 (Br. G.-O. 2, S. 508, 511, zum Teil abgedruckt W. VIII, S. 444—445).
\par
Die Abhandlung selbst ist 1828 in den Denkschriften der Göttinger Societät erschienen:
\par\noindent (1)\quad Disquisitiones generales circa superficies curvas, auctore \textsc{Carolo Friderico Gauss}, Societati regiae
oblatae d. 8. Octob. 1827. Commentationes societatis regiae scientiarum Gottingensis recentiores, Commentationes
classis mathematicae. T. VI (ad annos 1823—1827), Gottingae 1828, S. 99—146.
\clearpage\SourceLeaf{451}\ArithmeticHead{120}{L}{STÄCKEL, GAUSS ALS GEOMETER.}
\noindent\hspace*{\parindent}Es gibt Sonderabzüge mit den Seitenzahlen 1 bis 50 und einer besonderen Titelseite, die den Vermerk:
Gottingae, Typis Dieterichianis, 1828 trägt.
\par
Der lateinische Text wurde in der fünften, von \textsc{Liouville} besorgten Ausgabe des Werkes: \textsc{G. Monge},
Application de l’analyse à la géométrie, Paris 1850, S. 505—546 abgedruckt unter dem Titel:
\par\noindent (2)\quad Recherches sur la théorie générale des surfaces courbes, par \textsc{M. C. F. Gauss}.
\par
Es folgen zwei Übersetzungen ins Französische:
\par\noindent (3)\quad Recherches générales sur les surfaces courbes par \textsc{M. Gauss}. Traduit du latin par \textsc{M. T[iburce]
A[badie]}, ancien élève de l’École polytechnique, Nouvelles annales de mathématiques, t. 11, Paris 1852,
S. 195—258.
\par\noindent (4)\quad Recherches générales sur les surfaces courbes, par \textsc{M. C. F. Gauss}, traduites en français, suivies de
notes et d’études sur divers points de la théorie des surfaces et sur certaines classes de courbes, par \textsc{M.
E. Roger}, Paris 1855.
\par
Nach \textsc{H. D. Thompson} (siehe Nr. 9) ist von (4) eine weitere Ausgabe Grenoble 1870, Paris 1871
erschienen.
\par\noindent (5)\quad Die Werke von \textsc{Carl Friedrich Gauss}, herausgegeben von der Königlichen Gesellschaft der Wissenschaften
in Göttingen bringen die Disq. gen. im vierten Bande, Göttingen 1873, S. 217—258; ein zweiter,
unveränderter Abdruck ist 1880 herausgekommen.
\par
Es gibt zwei Übersetzungen ins Deutsche. Die erste ist ein Teil des Werkes: \textsc{O Böklen}, Analytische
Geometrie des Raumes, zweite Auflage, Stuttgart 1884, dessen zweiter Teil den Doppeltitel führt:
\par\noindent (6)\quad Disquisitiones generales circa superficies curvas von \textsc{C. F. Gauss}, ins Deutsche übertragen, mit Anwendungen
und Zusätzen. Die \textsc{Fresnel}sche Wellenfläche.
\par
Die Übersetzung steht S. 197—232. Die erste Auflage, Stuttgart 1861, enthält die Übersetzung der
Disq. gen. noch nicht.
\par
Zweitens ist zu nennen:
\par\noindent (7)\quad Allgemeine Flächentheorie (Disquisitiones generales circa superficies curvas) von \textsc{Carl Friedrich
Gauss} (1827). Deutsch herausgegeben von \textsc{A. Wangerin}. Heft 5 von \textsc{Ostwalds} Klassikern der exakten
Wissenschaften, Leipzig 1889, 62 S.; zweite revidierte Auflage, Leipzig 1900, 64 S.
\par
In den Budapester Mathematisch-physikalischen Blättern hat \textsc{Nikolaus Szíjártó} eine Übersetzung
ins Magyarische veröffentlicht:
\par\noindent (8)\quad A felületek általános elmélete. Irta \textsc{Gauss Károly Frigyes}. Forditotta \textsc{Szíjártó Miklós}. Mathematikai
és physikai lapok, Band 6, Budapest 1897, S. 45—114.
\par
Eine Übersetzung ins Englische enthält das Buch:
\par\noindent (9)\quad \textsc{Karl Friedrich Gauss}, General investigations of curved surfaces of 1827 and 1825. Translated
with notes and a bibliography by \textsc{J. C. Morehead} and \textsc{A. M. Hiltebeitel}. The Princeton University
Library, 1902.
\par
Die Einleitung von \textsc{H. D. Thompson} gibt bibliographische Notizen. Es folgt S. 1—44 die Übersetzung
der Disq. gen. Beigegeben sind Übersetzungen der Selbstanzeige und der 1900 im achten Bande
der Werke aus dem Nachlass herausgegebenen Neuen allgemeinen Untersuchungen über die krummen
Flächen.
\par\bigskip\begin{center}\rule{.18\linewidth}{.3pt}\end{center}
\endgroup
\clearpage\SourceLeaf{452}\ArithmeticHead{121}{R}{SCHLUSSBEMERKUNG.}
\begin{center}\textbf{SCHLUSSBEMERKUNG.}\end{center}
\begingroup\small
\noindent\hspace*{\parindent}Kurze Zeit nach dem Abdruck der vorstehenden Abhandlung in den \textbf{\textit{Materialien für eine wissenschaftliche
Biographie von Gauss}}, wurde der Verfasser \textsc{Paul Stäckel} aus voller Schaffenskraft durch einen
jähen Tod der Wissenschaft entrissen. — Seine grossen Verdienste um die Weiterführung der \textsc{Gauss}ausgabe
vom VIII. Bande an und um die würdige Schilderung der Leistungen von \textsc{Gauss} auf den verschiedenen
Gebieten der Mathematik erfahren mit dem hier wiedergegebenen Aufsatz ihre Krönung. — Als es sich nun
darum handelte, \textsc{Stäckels} Abhandlung über \textit{Gauss als Geometer} in den Band X\textsubscript{2} der Werke einzufügen,
konnte nur ein in allem Wesentlichen unveränderter Wiederabdruck aus den \textit{Materialien} in Frage kommen.
Der Unterzeichnete hat im Vereine mit \textsc{Friedrich Engel} eine sorgfältige Durchsicht der Arbeit
des dahingegangenen Freundes vorgenommen, wobei sich nur an einigen Stellen geringfügige Änderungen
als erforderlich erwiesen haben. Es ist uns aber bekannt, dass \textsc{Stäckel} selbst die Absicht hatte, beim
Wiederabdruck seines Aufsatzes in den Werken einen Punkt näher zu erörtern, den völlig aufzuklären ihm
bei der Abfassung noch nicht gelungen war. Es handelt sich nämlich darum, welche Bedeutung \textsc{Gauss}
der Ausmessung des Dreiecks Brocken, Hohenhagen, Inselsberg in bezug auf die Frage beigelegt hat, ob
man die Euklidische oder eine nichteuklidische Geometrie als theoretische Grundlage für die Messungen
auf der Erde und am Himmel anzunehmen habe. — Da sich im Nachlasse \textsc{Stäckels} keine Aufzeichnung
über diese Frage gefunden hat, so müssen wir uns damit begnügen, diejenigen Stellen aus \textsc{Sartorius von
Waltershausens} Schrift \textit{Gauss zum Gedächtnis} hier wiederzugeben, die sich darauf beziehen.
\par\medskip\begin{center}\rule{.18\linewidth}{.3pt}\end{center}
\endgroup
\noindent\textsc{Sartorius}, S. 53.
\par
». . . Das Heliotrop fand sogleich bei der Hannöverschen Triangulation
seine volle Anwendung und das grosse Dreieck, vielleicht das grösste, welches
gemessen worden ist, nämlich zwischen dem Brocken, dem Inselsberg und dem
Hohenhagen, wurde mit Hilfe desselben so genau gemessen, dass die Summe
der drei Winkel nur um etwa zwei Zehnteile einer Sekunde sich von zwei
Rechten entfernt.«
\par\medskip\noindent\textsc{Sartorius}, S. 81.
\par
». . . Die Geometrie betrachtete \textsc{Gauss} nur als ein konsequentes Gebäude,
nachdem die Parallelentheorie als Axiom an der Spitze zugegeben sei; er sei
indes zur Überzeugung gelangt, dass dieser Satz nicht bewiesen werden könne,
doch wisse man aus der Erfahrung, z. B. aus den Winkeln des Dreiecks
Brocken, Hohenhagen, Inselsberg, dass er näherungsweise richtig sei. Wolle
man dagegen das genannte Axiom nicht zugeben, so folge daraus eine selbständige
Geometrie, die er gelegentlich ein Mal verfolgt und mit dem Namen
Antieuklidische Geometrie bezeichnet habe.«
\par\smallskip\begingroup\small
Hierzu ist noch die oben S. 33 abgedruckte Stelle aus \textsc{Gauss}’ Brief an \textsc{Taurinus} vom 8. November
1824 zu vergleichen.\hfill\textsc{Schlesinger}.
\par\medskip\begin{center}\rule{.23\linewidth}{.3pt}\end{center}
\par\smallskip\noindent X\textsubscript{2} Abh. 4.\hfill 16\par\endgroup
\clearpage\SourceLeaf{453}\NoSourceHead
\begin{center}\textbf{Inhaltsverzeichnis.}\end{center}
\begingroup\fontsize{8.5}{11.5}\selectfont
\setlength{\tabcolsep}{3pt}
\newcommand{\TP}[1]{{\fontsize{8}{10}\selectfont\textls[350]{#1}}}
\newcommand{\CT}[4]{#1 & #2\ #3 & #4\\}
\noindent\begin{tabular}{@{}r p{.80\linewidth} r@{}}
\CT{1.}{Einleitung}{\TP{.............................}}{Seite 3}
\multicolumn{3}{c}{\textls[100]{I. Die Grundlagen der Geometrie.}}\\[3pt]
\CT{2.}{Allgemeines über die Arbeitsweise von \textsc{Gauss}}{\TP{................}}{— 6}
\multicolumn{3}{@{}p{.98\linewidth}@{}}{\textit{A. Von den Anfängen der nichteuklidischen Geometrie bis zur Entdeckung der transzendenten Trigonometrie (1792—1817).}}\\
\CT{3.}{Einleitendes. Die Jugendzeit (1792—1795)}{\TP{.................}}{— 17}
\CT{4.}{Fortschritte in den Grundlagen der Geometrie (1795—1799)}{\TP{...........}}{— 19}
\CT{5.}{Schwanken und Zweifel (1799—1805)}{\TP{...................}}{— 25}
\CT{6.}{Die Entdeckung der transzendenten Trigonometrie (1805—1817)}{\TP{..........}}{— 28}
\multicolumn{3}{@{}p{.98\linewidth}@{}}{\textit{B. Der Ausbau der nichteuklidischen Geometrie (seit 1817).}}\\
\CT{7.}{Die Zeit der Geodäsie und der Flächentheorie; \textsc{Schweikart} und \textsc{Taurinus} (1817—1831)}{}{— 31}
\CT{8.}{Die weitere Entwicklung bei \textsc{Gauss}; \textsc{Johann Bolyai} und \textsc{Lobatschefskij} (1831—1846)}{}{— 35}
\CT{9.}{Nachwirkung der \textsc{Gauss}schen Gedanken}{\TP{..................}}{— 41}
\multicolumn{3}{@{}p{.98\linewidth}@{}}{\textit{C. Sonstige Beiträge zur Axiomatik.}}\\
\CT{10.}{Weitere Untersuchungen über die Grundlagen der Geometrie}{\TP{...........}}{— 42}
\multicolumn{3}{c}{\textls[100]{II. Geometria situs.}}\\[3pt]
\CT{11.}{Allgemeines über die Geometria situs bei \textsc{Gauss}}{\TP{...............}}{— 49}
\CT{12.}{Verknotungen und Verkettungen von Kurven}{\TP{.}\hspace{.8em}\TP{..............}}{— 50}
\CT{13.}{\textsc{Möbius}, \textsc{Listing}, \textsc{Riemann}}{\TP{................}\hspace{.8em}\TP{......}}{— 51}
\multicolumn{3}{c}{\textls[100]{III. Die komplexen Grössen in ihrer Beziehung zur Geometrie.}}\\[3pt]
\CT{14.}{Kreisteilung}{\TP{................}\hspace{.8em}\TP{...........}}{— 56}
\CT{15.}{Elliptische, im besonderen lemniskatische Funktionen}{\TP{.........}\hspace{.8em}\TP{...}}{— 60}
\CT{16.}{Existenz der Wurzeln algebraischer Gleichungen}{\TP{...}\hspace{.8em}\TP{...........}}{— 62}
\CT{17.}{Biquadratische Reste}{\TP{..............}\hspace{.8em}\TP{.}\hspace{.8em}\TP{........}}{— 63}
\CT{18.}{Benutzung der komplexen Grössen für geometrische Untersuchungen}{\TP{....}\hspace{.8em}\TP{...}}{— 64}
\CT{19.}{Weiterentwicklung der Lehre von den komplexen Grössen}{\TP{............}}{— 65}
\CT{20.}{Komplexe Grössen mit mehr als zwei Einheiten}{\TP{...............}}{— 67}
\multicolumn{3}{c}{\textls[100]{IV. Elementare und analytische Geometrie.}}\\[3pt]
\CT{21.}{Allgemeines}{\TP{............................}}{— 69}
\CT{22.}{Das Dreieck}{\TP{............................}}{— 70}
\end{tabular}
\clearpage\SourceLeaf{454}\ArithmeticHead{123}{R}{INHALTSVERZEICHNIS.}
\noindent\begin{tabular}{@{}r p{.80\linewidth} r@{}}
\CT{23.}{Das Viereck}{\TP{............................}}{Seite 71}
\CT{24.}{Die Vielecke}{\TP{............................}}{— 74}
\CT{25.}{Der Kreis und die Kugel}{\TP{........................}}{— 76}
\CT{26.}{Kegelschnitte und Flächen zweiter Ordnung}{\TP{.................}}{— 78}
\CT{27.}{Sphärische Trigonometrie}{\TP{........................}}{— 81}
\multicolumn{3}{c}{\textls[100]{V. Die allgemeine Lehre von den krummen Flächen.}}\\[3pt]
\CT{28.}{Entwicklung der Grundgedanken bis zum Jahre 1816}{\TP{.............}}{— 84}
\CT{29.}{Die Kopenhagener Preisschrift (1822).}{\TP{....................}}{— 88}
\CT{30.}{Vorarbeiten zu den Allgemeinen Untersuchungen über die krummen Flächen (1822—1825)}{}{— 94}
\CT{31.}{Die Entstehung des Disquisitiones generales circa superficies curvas (1826—1827)}{\TP{...}}{— 101}
\CT{32.}{Weitere Untersuchungen über krumme Flächen}{\TP{................}}{— 106}
\CT{33.}{Bedeutung und Wirkung der Disquisitiones generales}{\TP{.............}}{— 110}
\CT{34.}{Bibliographischer Anhang}{\TP{........................}}{— 119}
\CT{}{\textls[100]{Schlussbemerkung}}{\TP{...........................}}{— 121}
\end{tabular}
\par\bigskip\begin{center}\rule{.25\linewidth}{.3pt}\end{center}
\endgroup\clearpage\endgroup
